\subsection{A compressed side circuit for the complex network}
\label{sec:complex-side-circuit}
Keep the complex motif of Section~\ref{sec:independent-complex}: $h=25$,
$v=\binom{25}3=2300$, $N=v^3$, $m=h^3=15625$, scalars in
$\mathbb Z[i,1/2]$, labels in $F=\mathbb F_2^{25}$ with the dot product, the
$h+1$ central wires, and the gather and scatter maps. For a triple $S$ write
$t_S$ for its indicator. The original side wires add
\begin{equation}\label{eq:complex-side-map}
 \frac12\sum_{T\cap S=\varnothing}x_T-\frac12\sum_{|T\cap S|=2}x_T
\end{equation}
to $y_S$. With the central coefficient $(|S\cap T|-1)/2$ this gives
$y\gets y+x$. They use one wire per ordered neighbor pair, that is
$vz_{\rm c}=3693800$ side wires per invocation. We compute the same map with
shared sums.

All sums below are formal sums of inputs $x_T$ whose two summands have
disjoint supports, so a source contributes to a node exactly when a directed
path joins them. Equal supports are identified within each of the two
families below, and nodes that feed no injected piece are removed.

\paragraph{Disjoint sums.}
Order the points $0,\ldots,24$. For a list of nodes indexed by distinct
points and $k\le3$, the routine $\mathrm{vec}$ returns, for every set $E$ of
at most $k$ of these points, the sum of the nodes whose points avoid $E$. It
splits the list at $\lfloor n/2\rfloor$, recurses, and forms each output by one
addition of a left and a right output whose exclusion sets have total size at
most $k$. For a point list $X$ and weights $w(u,v)$ on its pairs, the routine
$\mathrm{pairs}$ returns the sums over pairs avoiding $E$, $|E|\le3$: split $X$
into halves $L,R$, recurse on both, apply $\mathrm{vec}$ to each cross row
$(w(u,r))_{r\in R}$, then apply $\mathrm{vec}$ down the columns of these
row outputs with the remaining exclusion budget, and add the three disjoint
regions. The routine $\mathrm{tri}$ does the same for triple sums: with halves
$L,R$ of $X$, put
\[
 P_r=\mathrm{pairs}\bigl(L,(u,v)\mapsto x_{uvr}\bigr),\qquad
 P'_u=\mathrm{pairs}\bigl(R,(a,b)\mapsto x_{uab}\bigr),
\]
apply $\mathrm{vec}$ over $r\in R$ to $P_r[E_L]$ and over $u\in L$ to
$P'_u[E_R]$, and add the four disjoint regions: triples in $L$, triples in $R$,
two points in $L$ with one in $R$, and one point in $L$ with two in $R$.

At the top level $L=\{0,\ldots,11\}$ and $R=\{12,\ldots,24\}$. For a target
$S$ with $S_L=S\cap L$ and $S_R=S\cap R$ the disjoint pieces are
\[
 \mathrm{tri}(L)[S_L],\qquad \mathrm{tri}(R)[S_R],
\]
the two $\mathrm{vec}$ outputs over the halves $R_1=\{12,\ldots,17\}$ and
$R_2=\{18,\ldots,24\}$ of the family $P_r[S_L]$, and the two $\mathrm{vec}$
outputs over the halves $L_1=\{0,\ldots,5\}$ and $L_2=\{6,\ldots,11\}$ of
the family $P'_u[S_R]$. These six pieces partition the triples disjoint from
$S$. They are internal nodes of the full computations because
$\mathrm{vec}$ splits at the same halves.

\paragraph{Pair stars.}
For a pair $a<b$, list the other points as $u_1<\cdots<u_{23}$ and form
prefix sums $\pi_i=\sum_{j\le i}x_{abu_j}$ and suffix sums
$\sigma_i=\sum_{j\ge i}x_{abu_j}$. Target $\{a,b,u_i\}$ receives
$\pi_{i-1}$ and $\sigma_{i+1}$. Over its three pairs, $S$ receives exactly the
triples meeting it in two points.

\paragraph{Injection rule.}
Inject every disjoint piece with coefficient $1/2$ and every pair-star piece
with coefficient $-1/2$, subject to one rule: if the triples of a piece cover
every point outside $S$, inject its two summands instead, recursively. A single
triple covers three of the $22$ outside points, so the rule terminates. At
$h=25$ it never splits a disjoint piece. It splits $\pi_{22}$ for $i=23$ and
$\sigma_2$ for $i=1$. Each target therefore receives twelve pieces.

The circuit has $80595$ additions ($68595$ disjoint, $12000$ pair-star) and
$27600$ injected pieces. Exact support comparison checks every coefficient
of \eqref{eq:complex-side-map}: all $3693800$ nonzero entries, their signs,
and all required zeros.

\subsection{The transparent schedule}
Allocate reversible roles as for the bit network: a node's outgoing edges are
its later uses and its injected pieces; an input slot fans out, and an
addition reuses its first incoming role as a pivot and retires the second.
This uses $R=c+q=80595+27600=108195$ roles. Let $L$ be the product of the
node mixers, $V$ the copy into input slots, and $J$ the injection, which adds
$\pm\frac12$ times a piece role to its target. With $G,\mathcal R$ the
original gather and scatter, a forward invocation performs
\[
 L,\ -J,\ L^{-1},\ -\mathcal R,\ V,\ G,\ \mathcal R,\ L,\ J,\ L^{-1},\
 G^{-1},\ V^{-1}.
\]
For initial side values $z$ and central values $w$, the target changes by
\[
 -JLz-\mathcal Rw+\mathcal R(w+Gx)+JL(z+Vx)=(JLV+\mathcal RG)x=x,
\]
and the last four operations restore $z$ and $w$. The inverse schedule, with
logical source $Y$ and target $X$, gives $X\gets X-Y$; forward, inverse and
forward invocations send $(X,Y)$ to $(-Y,X)$ and restore every auxiliary
value. Every gate is pointwise linear with coefficients in
$\{0,\pm1,\pm\frac12\}$.

\subsection{Binary labels}
Assign each node $z$ a subspace $U_z\subset F$:
\begin{itemize}
\item an input node $T$ receives $\langle t_T\rangle$;
\item a disjoint-sum node receives the coordinate space
 $F_{X}=\langle e_x:x\in X\rangle$, where $X$ is the union of its triples;
\item a pair-star node for $\{a,b\}$ with leaf set $\Lambda$ receives
 $\langle t_{abu}:u\in\Lambda\rangle$.
\end{itemize}
The vectors $t_{abu}$ have weight three and pairwise intersections of size
two, so they form an orthonormal basis of the pair-star label. Coordinate
spaces have orthonormal coordinate bases. Every label is therefore
nondegenerate. Along a dataflow edge $z\to w$ the support grows, so
$U_z\subset U_w$. The injection gate for target $S$ uses the physical $Y$
frame, whose active component is $t_S^\perp$. Every triple of a piece meets
$S$ in an even number of points, so the piece label lies in $t_S^\perp$.

\begin{lemma}\label{lem:complex-piece-residual}
If a piece leaves some point $y\notin S$ uncovered, then the residual
$t_S^\perp\cap U_z^\perp$ of its injection edge has an orthonormal basis.
If instead a disjoint piece covers every point outside $S=\{a,b,c\}$, its
injection residual is $\langle e_a+e_b,e_b+e_c\rangle$, which has none.
\end{lemma}
\begin{proof}
Both $U_z$ and $t_S^\perp$ are nondegenerate, so the residual is
nondegenerate. The vector $e_y$ is orthogonal to $t_S$, to every coordinate
vector in a disjoint label, to every $t_{abu}$ with $u\ne y$, and to an
input line $t_T$ with $y\notin T$. Thus the residual contains a vector of norm
one and is nonalternating; the manuscript's absorption argument then gives
an orthonormal basis. In the second case
$t_S^\perp=F_{[h]\setminus S}\mathbin\perp\langle e_a+e_b,e_b+e_c\rangle$.
The displayed plane is nondegenerate, and all its vectors have even weight,
hence norm zero.
\end{proof}

A nondegenerate label between $F_{[h]\setminus S}$ and $t_S^\perp$ must be
one of these two spaces, since a proper nonzero part of the plane is an
isotropic line. A full-coverage disjoint piece thus cannot be repaired by
relabeling; the injection rule avoids it. Any partition of the disjoint
triples into pieces that each miss an outside point needs at least four
missed points, so four pieces per target are necessary for this label class.

The remaining residuals are covered by the following list, which the checker
verifies by exact binary linear algebra for every node, edge, piece and
compiled role.
\begin{itemize}
\item Coordinate to coordinate: $F_{X'\setminus X}$, a coordinate space.
\item Input line to coordinate space: $F_X\cap t_T^\perp$ contains $e_x$ for
 $x\in X\setminus T$, which is nonempty for an addition node.
\item Pair-star chain: the span of the new $t_{abu}$.
\item Entry from the early frame: the label itself.
\item Exit to the late frame: $U_z^\perp$, which contains a unit coordinate
 vector outside the covered points; for a pair-star node, outside
 $\{a,b\}\cup\Lambda$, which omits at least the excluded third point.
\end{itemize}
Each nonzero residual is nondegenerate and contains a vector of odd weight,
so it has an orthonormal basis.

\subsection{Frames in both stage directions}
Use the frames of the bit construction with these binary labels. With
$A=B\mathbin\perp P$ and future line factor $Q$, put $D_0=B\otimes F$,
$D_1=A\otimes F$ and $D_U=(B\otimes F)\mathbin\perp(P\otimes U)$, all
tensored by $Q$. Early mixers and the early injection use $D_0$; middle node
gates use $D_{U_z}$; copy gates use the physical $X$ frame
$D_{\langle t_T\rangle}$ and injection gates the physical $Y$ frame
$D_{t_S^\perp}$; late cleanup uses $D_1$. An edge residual is
$P\otimes E\otimes Q$ for a residual $E\subset F$ from the list above. The
lines $P,Q$ are spanned by tensor products of triple indicators, which have
norm one, so tensoring an orthonormal basis of $E$ gives one of the residual.
The data wires pass through the same frames as in the original motif and
have the same residuals.

In the inverse schedule the nontrivial labels occur in the first $L^{-1}$.
Assign $D_{U_z^\perp}$ to node $z$ there, $D_0$ to the preceding $L$, and $D_1$
to the second $L$, $J$ and $L^{-1}$. A forward inclusion $U_z\subset U_w$ gives
$U_w^\perp\subset U_z^\perp$ with the same residual $U_w\cap U_z^\perp$. A piece
role enters from the physical $X$ line $\langle t_S\rangle\subset U_z^\perp$
with residual $t_S^\perp\cap U_z^\perp$, the forward injection residual. An input
role starts at the physical $Y$ frame $B\otimes\langle t_T\rangle$ of the first
copy, passes $D_0$ with residual $B\otimes t_T^\perp$ as a data wire does, and
ends at $t_T^\perp$, exactly its physical $Y$ frame. Entries and exits
exchange the forward residuals $U_z$ and $U_z^\perp$. Thus every residual of the
inverse stage is one already shown to have an orthonormal basis.

\subsection{Rank, phases and depth}
\begin{proposition}[Compressed complex interface]\label{prop:compressed-complex-interface}
The complex network with this side circuit has
\begin{align*}
 W_{\rm c}&=2N+3v^2(R+h+1)=1741801270000,\\
 L_{\rm c}&=3v^2h(h+1)=10315500000,\\
 s_{\rm c}&=W_{\rm c}m-2N+2L_{\rm c}=27215641140750000,
\end{align*}
satisfies the binary phase interface of
Proposition~\ref{prop:compact-complex-interface} with these constants, and has
$\eta_{\rm c}=(W_{\rm c}m-s_{\rm c})/(W_{\rm c}m)=28/205789375$.
\end{proposition}
\begin{proof}
The schedule gives the signed exchange and restores arbitrary scratch. Every
side role has source label $0$ and sink label $F^{\otimes3}$, and its labels
increase, so it contributes exactly $m$ to the residual count. Data wires and
central wires are unchanged, so the only decreases remain the $h$-dimensional
central returns. Telescoping gives $s_{\rm c}$. Every residual has an
orthonormal basis, so each edge is a product of one translation kernel per
basis vector, as in the original phase interface. The endpoint phases of the
data roles and the signed-exchange correction are unchanged, and every
auxiliary role still runs from frame $I$ to $C^{\otimes m}$.

The coefficient-depth argument needs a bound on gates, not on gates per wire.
An invocation has $4\cdot82895$ mixer gates, $2v$ copy gates, $2v$ injection
gates (one per target) and four central gates. Over $3v^2$ invocations this is
$5408242080000\le12W_{\rm c}$. Each gate has at most $W_{\rm c}$ inputs and
outputs and coefficients in $\{0,\pm1,\pm\frac12\}$, so
$E=64(W_{\rm c}+m+1)^3$ bounds the depth of one invocation as before. Also
$2\le s_{\rm c}<m^5$, so the generalized stopped guard applies with these
constants. All gates, labels and bases are fixed finite data.
\end{proof}

The rational log bound $\log m<966/100$ and
$\eta_{\rm c}>(14/10^9)(966/100)$ give, by $e^{-x}>1-x$,
\begin{equation}\label{eq:compressed-complex-exponent}
 s_{\rm c}/W_{\rm c}<m^{\sigma},\qquad \sigma=1-\frac{14}{10^9}.
\end{equation}
This is $33.49$ times the previous complex saving $418/10^{12}$.
